\sloppy
\PassOptionsToPackage{table}{xcolor}

% ============================================================
%  preamble.inc.tex — Пакеты и настройки стиля ГОСТ 7-32
%  Настраиваемые параметры вынесены в config.inc.tex
% ============================================================

% --- Гиперссылки в PDF ---
\usepackage{float}
\usepackage[
  bookmarks=true, colorlinks=true, unicode=true,
  urlcolor=black, linkcolor=black, anchorcolor=black,
  citecolor=black, menucolor=black, filecolor=black,
]{hyperref}
\AfterHyperrefFix

% --- Подписи к рисункам / таблицам ---
\usepackage{ragged2e}
\usepackage{caption}
\captionsetup{%
   justification=raggedright,
   labelfont=bf,
   singlelinecheck=off,
   font=small,
   labelsep=endash
}
\captionsetup[figure]{font=small}
\captionsetup[table]{
   justification=raggedleft,
   labelfont=bf,
   singlelinecheck=off,
   font=small
}

% --- Типографика ---
\usepackage{microtype}
\ifInvisibleDashes
  \MakeDashesBold
\fi

% --- Графика ---
\usepackage{graphicx}
\usepackage{tikz}
\usetikzlibrary{arrows, positioning, shadows}
\usetikzlibrary{automata, arrows.meta, positioning, fit}

% --- Геометрия (значения по умолчанию, перезаписываются в config.inc.tex) ---
\geometry{ignorefoot}

% --- Списки и таблицы ---
\usepackage{enumerate}
\usepackage{multirow}
\usepackage{paralist,array}
\usepackage{longtable}
\usepackage{pdflscape}

% --- Рамка вокруг всех иллюстраций ---
% Каждое изображение обводится тонкой рамкой и вписывается в ширину строки
% с учётом самой рамки, поэтому width=\textwidth не вылезает на поля.
\usepackage{adjustbox}
\let\origincludegraphics\includegraphics
\renewcommand{\includegraphics}[2][]{%
    \adjustbox{cfbox=black!45 0.4pt 3pt, max width=\linewidth}{\origincludegraphics[#1]{#2}}%
}
% Уменьшенный шрифт во всех таблицах (документ набран 14pt).
\usepackage{etoolbox}
\AtBeginEnvironment{tabular}{\small}
\AtBeginEnvironment{longtable}{\footnotesize}
\usepackage{booktabs}

% --- Цветовое кодирование состояния в таблицах ---
\usepackage{xcolor}
\definecolor{cdone}{RGB}{223,240,216}   % реализовано
\definecolor{cpart}{RGB}{252,243,207}   % частично
\definecolor{cplan}{RGB}{219,233,248}   % целевое / не реализовано
\definecolor{cna}{RGB}{238,238,238}     % неприменимо
\definecolor{crisk}{RGB}{250,222,222}   % первоочередной риск
\newcommand{\sdone}{\cellcolor{cdone}}
\newcommand{\spart}{\cellcolor{cpart}}
\newcommand{\splan}{\cellcolor{cplan}}
\newcommand{\sna}{\cellcolor{cna}}
\newcommand{\srisk}{\cellcolor{crisk}}

% --- Разное ---
\usepackage{blindtext}
% \usepackage{syntax}  % КОНФЛИКТ с minted: syntax переопределяет \let в \makeatletter, ломая .pygstyle
\usepackage{verbatim}
\usepackage{pdfpages}
\usepackage{setspace}
\usepackage{multicol}
\usepackage{calc}
\usepackage{ifthen}
\usepackage{algorithm2e}
\usepackage{amsmath}
\usepackage{fontspec}
\usepackage[normalem]{ulem}
\useunder{\uline}{\ul}{}

% --- Таблицы из CSV ---
\usepackage{pgfplotstable}
\usepackage{csvsimple-l3}

% --- Листинги (minted) ---
\usepackage[newfloat]{minted}

\DeclareFloatingEnvironment[
  placement={!ht},
  name=Equation
]{eqndescNoIndent}
\edef\fixEqndesc{%
  \noexpand\setlength{\noexpand\parindent}{\the\parindent}%
  \noexpand\setlength{\noexpand\parskip}{\the\parskip}%
}
\newenvironment{eqndesc}[1][!ht]{%
    \begin{eqndescNoIndent}[#1]\fixEqndesc%
}{\end{eqndescNoIndent}}

\newenvironment{code}{\captionsetup{type=listing}}{}
\DeclareCaptionLabelFormat{myformat}{#1 #2}
\captionsetup[listing]{labelformat=myformat}
\renewcommand{\thelisting}{\textsc{\lowercase{\arabic{listing}}}}

\renewcommand\ListingInChapter{%
    \counterwithin{listing}{chapter}
    \renewcommand{\thelisting}{\textsc{\lowercase{\thechapter}}\arabic{listing}}
}

% --- Ширина бокса под номер главы в оглавлении ---
% В G2-105.sty задано 1em, чего хватает только на однозначные номера: при
% двузначном номере цифры наезжают на название главы. 1.6em согласуется с
% настройкой \l@part в том же файле и безопасно вмещает любые 2-значные номера.
\makeatletter
\renewcommand{\l@chapter}{\@dottedtocline{1}{0mm}{1.6em}}
\makeatother
